建设投入何时形成可交付能力，储能和协调如何缓解连续冲击，是多能源转型的关键问题。本文基于2010—2024年统计资料，构建政府、消费载能群、发电与储能投资、燃料供给四类主体，以217维状态、建设队列及三燃料守恒连接年度演化与小时运行，并在MATLAB/Simulink和AnyMath中完成实际仿真。

针对现状评价，2024年一次能源消费59.60亿吨标准煤，煤炭占53.2\%，风光发电占18.3\%；供给暴露、产区集中与燃料价格共同确定协同需求。针对双碳推演，基准2030年能源CO$_2$为11.351 Gt，至2060年窗口末端仍未确认达峰，剩余净排12.639 Gt；政策强化降至9.325 Gt。全因子机制对照表明，三类主动反馈合计减少终点能源CO$_2$ 0.283 Gt，同时增加累计未供有用服务3.011 PWh。

针对安全优化，以逐年五类服务缺失不超过1\%筛查2460组投资控制。名义折中方案在48组开发扰动中仅3组年度可靠，保供修复提高至20组，另48组独立设计中21组达标；名义折现资本增加26.526万亿元、2060年净排增加2.037 Gt。196条建设时点路径显示，窗口内扩建改善6组2030年油缺口，全期扩建却使可靠数由20降至19组，必须同时评价规模、资金配置与投产时钟。

小时层连续反馈实际库存与电量。协调相对可达全电源余量规则少缺供109.455 GWh（2.073\%），共同满电终点下收益保留；增加50\%储能功率在29/32格削峰，仅5格减少累计缺供。八组离线参照给出对应电量边界，四夜守恒解释原相位5.170 TWh下界。在该边界内，保峰仍可省费12.861—20.924亿元；再守排放后，六组只余698.316—2738.648万元，两组数值等费。建设兑现、时序分配、可交付能量与资源准备应据此分别配置。
